\uuid{FblR}
\titre{Matrice d'une application linéaire/Polynômes}
\auteur{hanani ; mzari}
\organisation{exo7}
\licence{CC BY-NC-SA 4.0}
\datecreate{2020-01-20}
\niveau{L1}
\module{Algèbre}
\chapitre{Matrice}
\sousChapitre{Matrice et application linéaire}
\difficulte{2}
\theme{Matrice d'une application linéaire, Polynômes}
\exo7qcmid{lille-2.matrice-application.013}

\contenu{
\begin{question}
\qclassification{108.03}
\qtags{motcle=Matrice d'une application linéaire/Polynômes}
Soit $\Rr_2[X]$ l'ensemble des polynômes à coefficients réels de degré $\le 2$, muni de sa base canonique ${\cal {B}} = \{1,X,X^2\}$ et $f$ l'application linéaire définie par :
$$\begin{array}{rccc}f:&\Rr_2[X] &\to& \Rr_2[X]\\
& P&\to &XP', \end{array}$$
où $P'$ est la dérivée de $P$. Soit ${\cal {B'}} = \{P_1,P_2,P_3\}$, où $P_1=1+X, P_2=1-X, P_3=(1+X)^2$. Quelles sont les assertions vraies ?
\begin{answers}
\good{${\cal {B'}}$ est une base de  $\Rr_2[X]$.}
\good{La matrice de $f$ dans la base $ {\cal {B}}$ est : $
\left(\begin{array}{rcc}
0&0&0\\
0&1&0\\
0&0&2\\
\end{array}\right)$}.
\good{La matrice de $f$ de la base $ {\cal {B'}}$ dans la base $ {\cal {B}}$ est : $
\left(\begin{array}{rcc}
0&0&0\\
1&-1&2\\
0&0&2\\
\end{array}\right)$}.
\bad{La matrice de $f$ dans la base $ {\cal {B'}}$ est : $
 \frac{1}{2}\left(\begin{array}{rcc}
1&-1&-4\\
-1&1&4\\
0&0&4\\
\end{array}\right)$}.
\end{answers}
\vskip2mm
\begin{explanations} La matrice de $f$ d'une base ${\cal {B}}=(u_j)$ dans une  base $ {\cal {B}}'=(v_i)$ est la matrice  dont la $j$ième colonne est constituée des coordonnées de  $f(u_j)$ dans la base $ {\cal {B}}'$. La matrice de $f$ dans la base $ {\cal {B'}}$ est : $
\frac{1}{2}\left(\begin{array}{rcc}
1&-1&-4\\
-1&1&0\\
0&0&4\\
\end{array}\right).$
\end{explanations}
\end{question}
}
